\documentclass[12pt]{article}
\usepackage[utf8]{inputenc}
\usepackage[T1]{fontenc}
\usepackage{hyperref}
\title{Zakhring Language Resources for K-12 Curriculum Development}
\author{Compiled by Claude Sonnet 5.5 (claude-sonnet-5-5) for Andr\'as Kornai}
\date{}
\begin{document}
\maketitle

\section{Introduction}
Zakhring (ISO 639-3 zkr; Glottocode zakh1243) is also called Meyor in India and Zha in China. Sources place its speakers in the Anjaw District of Arunachal Pradesh, along the Lohit River, and in a few villages in Tibet. Estimates range from about 900 to about 1,000 speakers. UNESCO lists it as definitely endangered. Its classification is disputed. It is variously treated as Sino-Tibetan, as East Bodish, as part of a Miju group, or possibly as an isolate. Little educational material exists, and I found none for schools.

\section{Primary Educational Resources}
I found no primers, textbooks, or school programmes for Zakhring.

\section{Linguistic and Reference Materials}
\begin{itemize}
\item Wikipedia, Zakhring language: an overview of classification, location, and speaker numbers. \url{https://en.wikipedia.org/wiki/Zakhring_language}
\item Wikipedia, Miju languages: the proposed grouping with Kaman. \url{https://en.wikipedia.org/wiki/Miju_languages}
\item Wikipedia, Zekhring people. \url{https://en.wikipedia.org/wiki/Zekhring_people}
\item Endangered Languages Project bibliography for Zakhring. \url{https://www.endangeredlanguages.com/lang/3140/bibliography}
\item A paper in \emph{Space and Culture, India} that discusses the community (not checked in detail). \url{https://www.spaceandculture.in/index.php/spaceandculture/article/download/152/89/}
\end{itemize}

\section{Religious and Community Resources}
\begin{itemize}
\item Asia Harvest people-group profile for the Zakhring, a missionary-oriented demographic profile. \url{https://people-groups.asiaharvest.org/India/Zakhring.pdf}
\end{itemize}
I found no scripture translations or community organizations.

\section{Digital and Audio Resources}
I found no audio, video, or software resources for Zakhring.

\section{Curriculum Development Context}
The language appears to be taught nowhere. A curriculum would need new documentation work done with the community, starting from the sources in the Endangered Languages Project bibliography. Because the language is spoken mainly in India, the relevant educational context is Arunachal Pradesh as well as China.
\end{document}